% ============ §19 Regular rings (原书 pp.153-161) ============
\section{Regular rings}\label{sec:19}\index{regular local rings, section 19}\index{regular rings, section 19}\index{minimal resolution, section 19}\index{global dimension, section 19}

\textbf{Minimal free resolutions.} Let $(A, m, k)$ be a local ring, $M$ and $N$ finite
$A$-modules. An $A$-linear map $\varphi\colon M \to N$ induces a $k$-linear map
$\bar{\varphi}\colon M \otimes k \to N \otimes k$, which we denote $\bar{\varphi}$; then one sees easily that
\[
  \varphi \text{ is an isomorphism } \iff \bar{\varphi} \text{ is surjective and } \Ker \varphi = mM.
\]
In particular for free modules $M$ and $N$, if $\bar{\varphi}$ is an isomorphism then $\rank
M = \rank N$, and writing $\varphi$ as a matrix we have $\det \varphi \notin m$, so that
\[
  \varphi \text{ is an isomorphism } \iff \bar{\varphi} \text{ is an isomorphism}.
\]
Let $M$ be a finite $A$-module. An exact sequence
\[
  (*) \quad \cdots \xrightarrow{d_{i+1}} L_i \xrightarrow{d_i} L_{i-1} \to \cdots \to L_1 \xrightarrow{d_1} L_0 \xrightarrow{\varepsilon} M \to 0
\]
(or the complex $L_\bullet$) is called a \emph{minimal (free) resolution}\index{minimal resolution} of $M$ if it satisfies the
three conditions (1) each $L_i$ is a finite free $A$-module, (2) $\bar{d}_i = 0$, or in other
words $d_iL_i \subset mL_{i-1}$ for all $i$, and (3) $\bar{\varepsilon}\colon L_0 \otimes k \to M \otimes k$ is an isomorph-
ism. Breaking up $(*)$ into short exact sequences $0 \to K_1 \to L_0 \to M \to 0$,
$0 \to K_2 \to L_1 \to K_1 \to 0, \ldots$, we have $L_0 \otimes k \cong M \otimes k$,
$L_1 \otimes k \cong K_1 \otimes k, \ldots$. Any two minimal resolutions of $M$ are isomorphic as com-
plexes (prove this!).

\begin{example*}
Let $x_1, \ldots, x_n \in m$ be an $A$-sequence, and let $K_\bullet = K_\bullet(x_1, \ldots, x_n)$
be the Koszul complex
\[
  0 \to K_n \to K_{n-1} \to \cdots \to K_0 \to A/(x_1, \ldots, x_n) \to 0;
\]
then $K_\bullet$ is a minimal resolution of $A/(x_1, \ldots, x_n)$ over $A$.
\end{example*}

Let $(A, m, k)$ be a Noetherian local ring; then a finite $A$-module $M$
always has a minimal resolution. Construction: let $\{\omega_1, \ldots, \omega_p\}$ be a
minimal basis of $M$, let $L_0 = Ae_1 + \cdots + Ae_p$ be a free module, and define
$\varepsilon\colon L_0 \to M$ by $\varepsilon(e_i) = \omega_i$; taking $K_1$ to be the kernel of $\varepsilon$ we get
$0 \to K_1 \to L_0 \to M \to 0$ with $L_0 \otimes k \cong M \otimes k$. Now $K_1$ is again a
finite $A$-module, so that we need only proceed as before.

\begin{lemma}\label{lem:19.1}
Let $(A, m, k)$ be a local ring, and $M$ a finite $A$-module. Suppose
that $L_\bullet$ is a minimal resolution of $M$; then
\begin{enumerate}[(i)]
  \item $\dim_k \Tor^A_i(M, k) = \rank L_i$ for all $i$,
  \item $\projdim M = \sup \{i \mid \Tor^A_i(M, k) \neq 0\} \le \projdim_A k$,
  \item if $M \neq 0$ and $\projdim M = r < \infty$ then for any finite $A$-module $N \neq 0$
        we have $\Ext^r_A(M, N) \neq 0$.
\end{enumerate}
\end{lemma}

\begin{proof}
(i) We have $\Tor^A_i(M, k) = H_i(L_\bullet \otimes k)$, but from the definition of
minimal resolution, $\bar{d}_i = 0$, and hence $H_i(L_\bullet \otimes k) = L_i \otimes k$, and the dimen-
sion of this as a $k$-vector space is equal to $\rank_A L_i$.

(ii) follows from (i).

(iii) Since $L_{r+1} = 0$ and $L_r \neq 0$, $\Ext^r_A(M, N)$ is the cokernel of $d_r^*\colon
\Hom_A(L_r, N) \to \Hom_A(L_{r-1}, N)$, but since $L_i$ is free, $\Hom_A(L_i, N)$ is just
a direct sum of a number of copies of $N$; we can write $d_r\colon L_r \to L_{r-1}$ as a
matrix with entries in $m$, and then $d_r^*$ is given by the same matrix, so that
$\Im(d_r^*) \subset m\,\Hom_A(L_r, N)$, and by NAK $\Ext^r_A(M, N) \neq 0$.
\end{proof}

\begin{remark}
One sees from the above lemma that $\Tor_i(M, k) = 0$ implies that
$L_i = 0$, and therefore $\projdim M < i$, so that $\Tor_j(M, k) = 0$ for $j > i$. It is
conjectured that this holds in more generality, or more precisely:
\emph{Rigidity conjecture.}\index{rigidity conjecture, section 19} Let $R$ be a Noetherian ring, $M$ and $N$ finite $R$-modules;
suppose that $\projdim M < \infty$. Then $\Tor^R_i(M, N) = 0$ implies that
$\Tor^R_j(M, N) = 0$ for all $j > i$.
\end{remark}

This has been proved by Lichtenbaum~\cite{Li1} if $R$ is a regular ring, but
is unsolved in general.

The following theorem is not an application of Lemma~\hyperref[lem:19.1]{1}, but is proved
by a similar technique.

\begin{theorem}[Auslander and Buchsbaum]\label{thm:19.1}
Let $A$ be a Noetherian local
ring and $M \neq 0$ a finite $A$-module. Suppose that $\projdim M < \infty$; then
\[
  \projdim M + \depth M = \depth A.
\]
\end{theorem}

\begin{proof}
Set $\projdim M = h$; we work by induction on $h$. If $h = 0$ then $M$
is a free $A$-module, so that the assertion is trivial. If $h = 1$, let
\[
  (\varphi) \quad 0 \to A^m \xrightarrow{\varphi} A^n \to M \to 0
\]
be a minimal resolution of $M$. We can write $\varphi$ as an $m \times n$ matrix with
entries in $m$. From $(\varphi)$ we obtain the long exact sequence
\[
  \cdots \to \Ext^i_A(k, A^n) \xrightarrow{\varphi^*} \Ext^i_A(k, A^m) \to \Ext^{i+1}_A(k, M) \to \cdots,
\]
and writing out $\Ext^i_A(k, A^m) = \Ext^i_A(k, A)^m$ and $\Ext^i_A(k, A^n) = \Ext^i_A(k, A)^n$,
we can express $\varphi^*$ by the same matrix as $\varphi$. However, the entries of $\varphi$ are
elements of $m$, and therefore annihilate $\Ext^i_A(k, A)$, so that $\varphi^* = 0$, and we
have an exact sequence
\[
  0 \to \Ext^i_A(k, A)^n \to \Ext^{i+1}_A(k, M) \to \Ext^{i+1}_A(k, A)^m \to 0
\]
for every $i$. Since $\depth M = \inf \{i \mid \Ext^i_A(k, M) \neq 0\}$ we have $\depth M =
\depth A - 1$ and the theorem holds if $h = 1$. If $h > 1$ then taking any exact
sequence
\[
  0 \to M' \to A^n \to M \to 0,
\]
we have $\projdim M' = h - 1$, so that an easy induction completes the
proof.
\end{proof}

\begin{lemma}\label{lem:19.2}
Let $A$ be a ring and $n > 0$ a given integer. Then the following
conditions are equivalent.
\begin{enumerate}[(1)]
  \item $\projdim M \le n$ for every $A$-module $M$;
  \item $\projdim M \le n$ for every finite $A$-module $M$;
  \item $\injdim N \le n$ for every $A$-module $N$;
  \item $\Ext^{n+1}_A(M, N) = 0$ for all $A$-modules $M$ and $N$.
\end{enumerate}
\end{lemma}

\begin{proof}
$(1) \implies (2)$ is trivial.

$(2) \implies (3)$ For any ideal $I$, the $A$-module $A/I$ is finite, so that
$\Ext^{n+1}_A(A/I, N) = 0$, and hence by \S\ref{sec:18}, Lemma~\hyperref[lem:18.1]{1}, $\injdim N \le n$.

$(3) \implies (4)$ is trivial, and $(4) \implies (1)$ is well-known (see p.~280).
\end{proof}

We define the \emph{global dimension}\index{global dimension} of a ring $A$ by
\[
  \gldim A = \sup \{\projdim M \mid M \text{ is an $A$-module}\}.
\]
According to Lemma~\hyperref[lem:19.2]{2} above, this is also equal to the maximum projective
dimension of all finite $A$-modules. If $(A, m, k)$ is a Noetherian local ring then by
Lemma~\hyperref[lem:19.1]{1}, $\gldim A = \projdim_A k$.

We have defined regular local rings (see \S\ref{sec:14}) as Noetherian local rings
for which $\dim A = \embdim A$, and we have seen that they are integral
domains (Theorem~\ref{thm:14.3}) and CM rings (Theorem~\ref{thm:17.8}). A regular local
ring is Gorenstein (Theorem~\ref{thm:18.1}, (5')). A necessary and sufficient condition
for a Noetherian local ring $(A, m, k)$ to be regular is that $\gr_m(A)$ is a poly-
nomial ring over $k$ (Theorem~\ref{thm:17.10}). The following theorem gives another
important necessary and sufficient condition.

\begin{theorem}[Serre]\label{thm:19.2}
Let $A$ be a Noetherian local ring. Then
\[
  A \text{ is regular} \iff \gldim A = \dim A \iff \gldim A < \infty.
\]
\end{theorem}

\begin{proof}
(I) Suppose that $(A, m, k)$ is an $n$-dimensional regular local ring.
Let $x_1, \ldots, x_n$ be a regular system of parameters; then since this is an
$A$-sequence, the Koszul complex $K_\bullet(x_1, \ldots, x_n)$ is a minimal free resolution of
$A/(x_1, \ldots, x_n) = k$, and $K_n \neq 0$, $K_{n+1} = 0$, so that as we have already seen,
$\gldim A = \projdim k = n$.

(II) Conversely, suppose that $\gldim A = r < \infty$ and $\embdim A = s$. We
prove that $A$ is regular by induction on $s$; we can assume that $s > 0$, that
is $m \neq 0$. Then $m \notin \Ass(A)$: for if $0 \neq a \in A$ is such that $ma = 0$, consider
a minimal resolution
\[
  0 \to L_r \to L_{r-1} \to \cdots \to L_0 \to k \to 0
\]
of $k$ (with $r > 0$); then $L_r \subset mL_{r-1}$, but then $aL_r = 0$, which contradicts the
assumption that $L_r$ is a free module. Thus we can choose $x \in m$ not
contained in $m^2$ or in any associated prime of $A$. Then $x$ is $A$-regular,
hence also $m$-regular, so that if we set $B = A/xA$ then according to Lemma~\hyperref[lem:18.2]{2}
of \S\ref{sec:18}, $\Ext^q_A(m, N) = \Ext^q_B(m/xm, N)$ for all $B$-modules $N$, and hence
we obtain $\projdim_B m/xm \le r$.

Now we prove that the natural map $m/xm \to m/xA$ splits, so that
$m/xA$ is isomorphic to a direct summand of $m/xm$. Since $x \notin m^2$, we can take
a minimal basis $x_1 = x, x_2, \ldots, x_s$ of $m$ starting with $x$ (here $s = \embdim
A$). We set $b = (x_2, \ldots, x_s)$, so that by the minimal basis condition,
$b \cap xA \subset xm$, and therefore there exists a chain
\[
  m/xA = (b + xA)/xA \cong b/(b \cap xA) \to m/xm \to m/xA
\]
of natural maps, whose composite is the identity. This proves the above
claim. Now clearly,
\[
  \projdim_B m/xA \le \projdim_B m/xm \le r.
\]
Taking a minimal $B$-projective resolution of $m/xA$ and patching it together
with the exact sequence $0 \to m/xA \to B \to k \to 0$ gives a projective
resolution of $k$ of length $\le r + 1$, and hence $\gldim B = \projdim_B k \le r + 1$,
so that by induction, $B$ is a regular local ring. Since $x$ is not contained
in any associated prime of $A$ we have $\dim B = \dim A - 1$, and therefore
$A$ is regular.
\end{proof}

\begin{theorem}[Serre]\label{thm:19.3}
Let $A$ be a regular local ring and $P$ a prime ideal; then
$A_P$ is a regular local ring.
\end{theorem}

\begin{proof}
Since $\projdim_A A/P \le \gldim A < \infty$, as an $A$-module $A/P$ has a
projective resolution $L_\bullet$ of finite length. Then $L_\bullet \otimes_A A_P$ is a projective
resolution of $(A/P) \otimes_A A_P = A_P/PA_P = K(P)$ as an $A_P$-module, so that $K(P)$
has a projective resolution of finite length as an $A_P$-module, which means that $A_P$
has finite global dimension; thus by the previous theorem, $A_P$ is
regular.
\end{proof}

\begin{definition*}
A \emph{regular ring}\index{regular rings} is a Noetherian ring such that the localisation
at every prime is a regular local ring. By the previous theorem, it is
enough\index{regular local rings} for the localisation at every maximal ideal to be regular.
\end{definition*}

\begin{theorem}\label{thm:19.4}
A regular ring is normal.
\end{theorem}

\begin{proof}
The definition of normal is local, so that it is enough to show that a
regular local ring is normal. We show that the conditions of the corollary
to Theorem~\ref{thm:11.5} are satisfied. (a) The localisation at a height 1 prime ideal is a
DVR by the previous theorem and Theorem~\ref{thm:11.2}. (b) All the prime
divisors of a non-zero principal ideal have height 1 by Theorem~\ref{thm:17.8} (the
implication regular $\Rightarrow$ CM).
\end{proof}

\begin{theorem}\label{thm:19.5}
If $A$ is regular then so are $A[X]$ and $A[[X]]$.
\end{theorem}

\begin{proof}
For $A[X]$, let $P$ be a maximal ideal of $A[X]$ and set $P \cap A = m$.
Then $A[X]_P$ is a localisation of $A_m[X]$, so that replacing $A$ by $A_m$, we can assume
$A$ is a regular local ring. Then setting $A/m = k$ we have $A[X]/m[X] =
k[X]$, so that there is a monic polynomial $f(X)$ with coefficients in $A$ such that
$P = (m, f(X))$, and such that $f$ reduces to an irreducible polynomial in
$k[X]$ modulo $m$. Then by Theorem~\ref{thm:15.1}, we clearly have
\[
  \dim A[X]_P = \htop P = 1 + \htop m = 1 + \dim A;
\]
on the other hand $m$ is generated by $\dim A$ elements, so that $P = (m, f)$ is
generated by $\dim A + 1$ elements, and therefore $A[X]_P$ is regular.

For $A[[X]]$, set $B = A[[X]]$ and let $M$ be a maximal ideal of $B$; then $X \in M$
by Theorem~\ref{thm:8.2}, (i). Therefore $M \cap A = m$ is a maximal ideal of $A$. Now
although we cannot say that $B_M$ contains $A_m[[X]]$, the two have the same
completion, $(B_M)^\wedge = (A_m)^\wedge[[X]]$. A Noetherian local ring is regular if and
only if its completion is regular (since both the dimension and embedding
dimension remain the same on taking the completion). Thus if we replace
$A$ by $(A_m)^\wedge$, the maximal ideal of $B = A[[X]]$ is $M = (m, X)$, and $\htop M =
\htop m + 1$, so that $B$ is also regular.
\end{proof}

Next we discuss the properties of modules which have finite free
resolutions (the definition is given below).

\begin{lemma}[Schanuel]\label{lem:19.3}
Let $A$ be a ring and $M$ an $A$-module. Suppose that
\[
  0 \to K \to P \to M \to 0 \quad \text{and} \quad 0 \to K' \to P' \to M \to 0
\]
are exact sequences with $P$ and $P'$ projective. Then $K \oplus P' \cong K' \oplus P$.
\end{lemma}

% 注: 原书此处用 λ, λ' 与矩阵 (φ, ψ, φψ = 1) 的写法, 扫描版矩阵不可辨;
% 此处按标准 fibre product 证法转录, 结论与原书一致.
\begin{proof}
From the fact that $P$ and $P'$ are projective, there exist $\lambda\colon P \to P'$
and $\lambda'\colon P' \to P$, giving the diagram
\[
  \begin{array}{ccccccc}
    0 & \to & K & \to & P & \xrightarrow{\varepsilon} & M \to 0 \\
    & & & & \downarrow \lambda' & & \downarrow \lambda \\
    0 & \to & K' & \to & P' & \xrightarrow{\varepsilon'} & M \to 0
  \end{array}
\]
with $\varepsilon'\lambda = \varepsilon$ and $\varepsilon\lambda' = \varepsilon'$.
We add in harmless summands $P'$ and $P$ to the two exact rows, and line
up the middle terms:
\[
  0 \to K \oplus P' \to P \oplus P' \xrightarrow{(\varepsilon, 0)} M \to 0
  \quad \text{and} \quad
  0 \to P \oplus K' \to P \oplus P' \xrightarrow{(0, \varepsilon')} M \to 0.
\]
Consider the submodule $W = \{(x, y) \in P \oplus P' \mid \varepsilon(x) = \varepsilon'(y)\}$ of $P \oplus P'$,
the fibre product of $P$ and $P'$ over $M$. Projection onto the first factor gives
the exact sequence $0 \to K' \to W \to P \to 0$: indeed $W \to P$ is surjective since
$\varepsilon'$ is surjective, and its kernel is $\{(0, y) \mid \varepsilon'(y) = 0\} = K'$. Since $P$ is
projective, this sequence splits, so that $W \cong K' \oplus P$. Similarly, projection onto
the second factor gives the exact sequence $0 \to K \to W \to P' \to 0$, which splits, so
that $W \cong K \oplus P'$. Therefore $\varphi$ induces an isomorphism
$K \oplus P' \cong P \oplus K'$.
\end{proof}

\begin{lemma}[generalised Schanuel lemma]\label{lem:19.4}
Let $A$, $M$ be as above, and suppose
that
\[
  0 \to P_n \to \cdots \to P_1 \to P_0 \to M \to 0
  \quad \text{and} \quad
  0 \to Q_n \to \cdots \to Q_1 \to Q_0 \to M \to 0
\]
are exact sequences with $P_i$ and $Q_i$ projective for
$0 \le i \le n - 1$. Then
\[
  P_0 \oplus Q_1 \oplus P_2 \oplus \cdots \cong Q_0 \oplus P_1 \oplus Q_2 \oplus \cdots
\]
(the alternating sums of $P_i$ and $Q_i$).
\end{lemma}

\begin{proof}
Write $K$ for the kernel of $P_0 \to M$ and $K'$ for the kernel of
$Q_0 \to M$; then, by the previous lemma, $K \oplus Q_0 \cong P_0 \oplus K'$. Now add
in harmless summands $Q_n$ and $P_n$ to $0 \to P_n \to \cdots \to P_1 \to
K \to 0$ and $0 \to Q_n \to \cdots \to Q_1 \to K' \to 0$ respectively, to obtain
exact sequences
\[
  0 \to P_n \to \cdots \to P_1 \oplus Q_0 \to K \oplus Q_0 \to 0
  \quad \text{and} \quad
  0 \to Q_n \to \cdots \to Q_1 \oplus P_0 \to P_0 \oplus K' \to 0,
\]
and induction on $n$ gives the assertion.
\end{proof}

A \emph{finite free resolution} (or \emph{FFR}\index{FFR} for short) of an $A$-module $M$ is a
resolution
\[
  0 \to F_r \to \cdots \to F_1 \to F_0 \to M \to 0 \quad \text{(of finite length)}
\]
with each $F_i$ a finite free module. If $M$ has an FFR we set
\[
  \chi(M) = \sum_{i=0}^{r} (-1)^i \rank F_i,
\]
and call $\chi(M)$ the \emph{Euler number}\index{Euler number} of $M$. By Lemma~\hyperref[lem:19.4]{4},
$\chi(M)$ is independent of the choice of FFR. Moreover, since for any prime
ideal $P$,
\[
  0 \to (F_r)_P \to \cdots \to (F_1)_P \to (F_0)_P \to M_P \to 0
\]
is an FFR of the $A_P$-module $M_P$ we have $\chi(M) = \chi(M_P)$. If $M$ is itself
free then one sees easily from Lemma~\hyperref[lem:19.4]{4} that $\chi(M) = \rank M$.

\begin{theorem}\label{thm:19.6}
Let $(A, m)$ be a local ring, and suppose that for any finite
subset $E \subset m$ there exists $0 \neq y \in A$ such that $yE = 0$; then the only
$A$-modules having an FFR are the free modules.
\end{theorem}

\begin{remark}
If $A$ is Noetherian then the assumption on $m$ is equivalent to
$m \in \Ass(A)$, or $\depth A = 0$. In this case the theorem is a special case of
Theorem~\ref{thm:19.1}.
\end{remark}

\begin{proof}
Suppose that $0 \to F_r \to F_{r-1} \to \cdots \to F_0 \to M \to 0$ is an FFR
of $M$, and set $N = \Coker(F_r \to F_{r-1})$; if we prove that $N$ is free then
we can decrease $r$ by 1, so that we only need consider the case $0 \to
F_1 \to F_0 \to M \to 0$. Now let $0 \to L_1 \to L_0 \to M \to 0$ be a minimal
free resolution of $M$; then since $L_0$ and $F_0$ are finitely generated, by
Schanuel's lemma (or by Theorem~\ref{thm:2.6}), $L_1$ is also finite. Considering bases of
$F_1$ and $L_1$, we can write down a set of generators of $L_1$ as a submodule of
$F_0$ using only a finite number of elements of $m$. Then by assumption, there
exists $0 \neq y \in A$ such that $yL_1 = 0$. Since $L_1$ is a free module, we must
have $L_1 = 0$, so that $M \cong L_0$, and is free.
\end{proof}

\begin{theorem}\label{thm:19.7}
Let $A$ be any ring; if $M$ is an $A$-module having an FFR then
$\chi(M) \ge 0$.
\end{theorem}

\begin{proof}
Choose a minimal prime ideal $P$ of $A$; since $\chi(M) = \chi(M_P)$, we can
replace $A$ by $A_P$, and then $A$ is a local ring with maximal ideal $m$ equal to
$\nil(A)$. Then the assumption of the previous theorem is satisfied: for
$x_1, \ldots, x_r \in m$, we can assume by induction on $r$ that there is a $z \neq 0$
such that $zx_1 = \cdots = zx_{r-1}$; but $x_r$ is nilpotent, so that there is an $i \ge 0$
such that $zx_r^i \neq 0$ but $zx_r^{i+1} = 0$, and we can take $y = zx_r^i$. Therefore by
the previous theorem $M$ is a free module, and $\chi(M) = \rank M \ge 0$.
\end{proof}

\begin{theorem}[Auslander and Buchsbaum~\cite{AB2}]\label{thm:19.8}
Let $A$ be a Noetherian
ring and $M$ an $A$-module, and suppose that $M$ has an FFR. Then the
following three conditions are equivalent:
\begin{enumerate}[(1)]
  \item $\ann(M) \neq 0$;
  \item $\chi(M) = 0$;
  \item $\ann(M)$ contains an $A$-regular element.
\end{enumerate}
\end{theorem}

\begin{proof}
$(1) \implies (2)$ Suppose that $\chi(M) > 0$; then for any $P \in \Ass(A)$ we have
$\chi(M_P) > 0$, and hence $M_P \neq 0$. By Theorem~\hyperref[thm:19.6]{6}, $M_P$ is a free $A_P$-module, so
that setting $I = \ann(M)$ we have $I_P = \ann(M_P) = 0$. If we set $J = \ann(I)$
then this is equivalent to $J \not\subset P$. Since this holds for every $P \in \Ass(A)$
we see that $J$ contains an $A$-regular element, but then $J \cdot I = 0$ implies that
$I = 0$, a contradiction.

$(2) \implies (3)$ If $\chi(M) = 0$ then by Theorem~\hyperref[thm:19.6]{6}, $M_P = 0$ for every $P \in \Ass(A)$.
This means that $\ann(M) \not\subset P$, so that $\ann(M)$ contains an $A$-regular
element.

$(3) \implies (1)$ is obvious.
\end{proof}

\begin{theorem}[Vasconcelos~\cite{Vas1}]\label{thm:19.9}
Let $A$ be a Noetherian local ring, and
$I$ a proper ideal of $A$; assume that $\projdim I < \infty$. Then
\[
  I \text{ is generated by an $A$-sequence} \iff I/I^2 \text{ is a free module over } A/I.
\]
\end{theorem}

\begin{proof}
$(\implies)$ is already known (Theorem~\ref{thm:16.2}). In fact, $I^v/I^{v+1}$ is a free
$A/I$-module for $v = 1, 2, \ldots$

$(\Longleftarrow)$ We can assume that $I \neq 0$. Since $I$ has finite projective dimension
over $A$ so has $A/I$, and since $A$ is local, $A/I$ has an FFR. Now $\ann(A/I) = I$,
so that by the previous theorem $I$ is not contained in any associated prime
of $A$, and therefore we can choose an element $x \in I$ such that $x$ is not
contained in $m^2$ or in any associated prime of $A$. Then $x$ is $A$-regular,
and $\bar{x} = x \bmod I^2$ is a member of a basis of $I/I^2$ over $A/I$; let $\bar{x}$,
$\bar{y}_2, \ldots, \bar{y}_n \in I$ be such that their images form a basis of $I/I^2$. Then if we
set $B = A/xA$, we see by the same argument as in (II) of the proof of
Theorem~\hyperref[thm:19.2]{2} that $\projdim_B I/xI < \infty$, and that $I/xA$ is isomorphic to a direct
summand of $I/xI$. We now set $I^* = I/xA$, so that $\projdim_A I^* < \infty$. But
on the other hand one sees easily that $I^*/I^{*2}$ is a free module over $B/I^*$,
and an induction on the number of generators of $I$ completes the proof.
\end{proof}

\begin{remark}
In Lech~\cite{Le1}, a set $x_1, \ldots, x_n$ of elements of $A$ is defined to be
\emph{independent}\index{independent elements} if
\[
  \sum a_ix_i = 0 \Rightarrow a_i \in (x_1, \ldots, x_n) \quad \text{for all } i.
\]
If we set $I = (x_1, \ldots, x_n)$ then this condition is equivalent to saying that
the images of $x_1, \ldots, x_n$ in $I/I^2$ form a basis of $I/I^2$ over $A/I$. Then if $A$
and $I$ satisfy the hypotheses of the previous theorem, the theorem tells us
that $I = (y_1, \ldots, y_n)$ with $y_1, \ldots, y_n$ an $A$-sequence. Setting $x_i = \sum a_{ij}y_j$ we
see that the matrix $(a_{ij})$ is invertible when considered in $A/I$; this means that
the determinant of $(a_{ij})$ is not in the maximal ideal of $A$, and so $(a_{ij})$ itself
is invertible. Thus $x_1, \ldots, x_n$ is an $A$-quasi-regular sequence, hence an
$A$-sequence. In particular, we get the following corollary.
\end{remark}

\begin{corollary*}
Let $(A, m)$ be a regular local ring. Then if $x_1, \ldots, x_n \in m$ are
independent in the sense of Lech, they form an $A$-sequence.
\end{corollary*}

However, if we try to prove this corollary as it stands, the induction
does not go through. The key to success with Vasconcelos' theorem is to
strengthen the statement so that induction can be used effectively. Now
as Kaplansky has also pointed out, the main part of Theorem~\hyperref[thm:19.2]{2} (the
implication $\gldim A < \infty \Rightarrow$ regular) follows at once from Theorem~\hyperref[thm:19.9]{9},
because if $m$ is generated by an $A$-sequence then $\embdim A \le \depth A \le
\dim A$.

% ============ 习题 ============
\section*{Exercises to \S\ref{sec:19}}
\emph{Prove the following propositions.}

\begin{exercise}\label{ex:19.1-ex}
Let $k$ be a field and $R = R_0 + R_1 + R_2 + \cdots$ a Noetherian graded ring
with $R_0 = k$; set $m = R_1 + R_2 + \cdots$. Show that if $R_m$ is an $n$-dimensional
regular local ring then $R$ is a polynomial ring $R = k[y_1, \ldots, y_n]$ with $y_i$
homogeneous of positive degree.
\end{exercise}

\begin{exercise}\label{ex:19.2-ex}
Let $A$ be a ring and $M$ an $A$-module. Say that $M$ is \emph{stably free}\index{stably free} if there exist
finite free modules $F$ and $F'$ such that $M \oplus F = F'$. Obviously a stably free
$A$-module $M$ is a finite projective $A$-module, and has an FFR $0 \to
F \to F' \to M \to 0$. Prove that, conversely, a finite projective module
having an FFR is stably free.
\end{exercise}

\begin{exercise}\label{ex:19.3-ex}
Prove that if every finite projective module over a Noetherian ring $A$ is
stably free then every finite $A$-module of finite projective dimension has an
FFR.
\end{exercise}

\begin{exercise}\label{ex:19.4-ex}
Prove that if every finite module over a Noetherian ring $A$ has an FFR
then $A$ is regular.
\end{exercise}
